% META: {"id": "1SPE-VARALEA-ALG-001", "chapitre": "1SPE-VARIABLES-ALEATOIRES", "type_objet": "algorithme", "capacites_codes": ["C2"], "programme_refs": ["PROB-VA-ALG-001"], "mode_creation": "ex_nihilo", "sources_inspiration": [], "status": "verified"}

\section{Algorithmique et expérimentations sur une variable aléatoire}

\subsection{Calculer l'espérance, la variance et l'écart type}

Une variable aléatoire finie prend ici les valeurs $-2$, $1$ et $4$ avec
les probabilités respectives $0{,}2$, $0{,}5$ et $0{,}3$. Le programme
suivant reçoit les valeurs et leurs probabilités; il ne dépend donc pas de
la situation choisie.

\begin{python}
from math import sqrt

def caracteristiques(valeurs, probabilites):
    """Renvoie (esperance, variance, ecart_type) d'une loi finie."""
    if len(valeurs) == 0 or len(valeurs) != len(probabilites):
        raise ValueError("la loi doit contenir autant de valeurs que de probabilites")
    if any(p < 0 for p in probabilites):
        raise ValueError("les probabilites doivent etre positives")
    if abs(sum(probabilites) - 1) > 1e-12:
        raise ValueError("la somme des probabilites doit valoir 1")
    esperance = sum(x * p for x, p in zip(valeurs, probabilites))
    variance = sum(
        (x - esperance) ** 2 * p
        for x, p in zip(valeurs, probabilites)
    )
    return esperance, variance, sqrt(variance)
\end{python}

\exemple{
Avec \code{valeurs = (-2, 1, 4)} et
\code{probabilites = (0.2, 0.5, 0.3)}, la fonction renvoie
\[
  E(X)=1{,}3,\qquad V(X)=4{,}41,\qquad \sigma(X)=2{,}1.
\]
Ces résultats sont des calculs exacts à l'arrondi d'affichage près : ils
proviennent de la loi donnée, sans simulation.
}

% META: {"id": "1SPE-VARALEA-ALG-002", "chapitre": "1SPE-VARIABLES-ALEATOIRES", "type_objet": "algorithme", "capacites_codes": ["C1"], "programme_refs": ["PROB-VA-ALG-002"], "mode_creation": "ex_nihilo", "sources_inspiration": [], "status": "verified"}
\subsection{Comparer les fréquences des lettres de deux textes}

Une lettre choisie au hasard dans un texte définit une variable aléatoire
finie. Sa loi est le tableau des fréquences des lettres. Pour rendre une
comparaison français--anglais pertinente, la fonction ci-dessous ignore la
casse, les espaces et la ponctuation; elle ramène aussi une lettre accentuée
à la lettre latine correspondante.

\begin{python}
from collections import Counter
from unicodedata import normalize

def frequences_lettres(texte):
    """Renvoie les frequences des lettres a...z presentes dans texte."""
    decompose = normalize("NFD", texte.lower())
    lettres = [caractere for caractere in decompose if "a" <= caractere <= "z"]
    if not lettres:
        return {}
    effectifs = Counter(lettres)
    total = len(lettres)
    return {
        lettre: effectifs[lettre] / total
        for lettre in sorted(effectifs)
    }
\end{python}

\propriete[Démarche de comparaison]{%
Appliquer la fonction à deux textes de longueurs comparables, puis comparer
les lettres les plus fréquentes. Les fréquences calculées décrivent les
textes choisis : un extrait différent, même dans la même langue, peut donner
un classement différent.
}

% BEGIN-VERIFY
% from collections import Counter
% from unicodedata import normalize
%
% def frequences_lettres(texte):
%     decompose = normalize("NFD", texte.lower())
%     lettres = [caractere for caractere in decompose if "a" <= caractere <= "z"]
%     if not lettres:
%         return {}
%     effectifs = Counter(lettres)
%     total = len(lettres)
%     return {lettre: effectifs[lettre] / total for lettre in sorted(effectifs)}
%
% # Une loi de probabilite : des frequences positives dont la somme vaut 1.
% loi = frequences_lettres("Il etait une fois")
% assert all(0 < f <= 1 for f in loi.values())
% assert abs(sum(loi.values()) - 1) < 1e-12
% # La casse, les espaces et la ponctuation sont ignores.
% assert frequences_lettres("Abc, abc !") == frequences_lettres("abcABC")
% # Une lettre accentuee est ramenee a la lettre latine correspondante.
% assert frequences_lettres("ete") == frequences_lettres("été")
% assert sorted(frequences_lettres("été")) == ["e", "t"]
% assert abs(frequences_lettres("été")["e"] - 2 / 3) < 1e-12
% # Un texte sans lettre ne definit aucune loi.
% assert frequences_lettres("123 ... !") == {}
% END-VERIFY

% META: {"id": "1SPE-VARALEA-EXP-001", "chapitre": "1SPE-VARIABLES-ALEATOIRES", "type_objet": "experimentation", "capacites_codes": ["C3"], "programme_refs": ["VA-EXP-SIMULER"], "mode_creation": "ex_nihilo", "sources_inspiration": [], "status": "verified"}

\subsection{Simuler une variable aléatoire finie}

On utilise la loi
\[
\begin{array}{c|ccc}
  x_i & -2 & 1 & 4\\ \hline
  P(X=x_i) & 0{,}2 & 0{,}5 & 0{,}3
\end{array}
\]
La fonction \code{tirage\_variable} partage l'intervalle $[0\,;1[$ en
trois intervalles de longueurs $0{,}2$, $0{,}5$ et $0{,}3$.

\begin{python}
from random import Random

VALEURS = (-2, 1, 4)
PROBABILITES = (0.2, 0.5, 0.3)

def tirage_variable(generateur):
    u = generateur.random()
    cumul = 0
    for valeur, probabilite in zip(VALEURS, PROBABILITES):
        cumul = cumul + probabilite
        if u < cumul:
            return valeur
    return VALEURS[-1]

def simuler_variable(n, graine):
    """Simule un echantillon de taille n de la variable X."""
    if n <= 0:
        raise ValueError("n doit etre strictement positif")
    generateur = Random(graine)
    return [tirage_variable(generateur) for _ in range(n)]
\end{python}

\erreurFrequente{
La graine sert à retrouver exactement le même échantillon lors d'une
vérification. La liste obtenue est une réalisation possible de la variable,
pas la loi elle-même.
}

% META: {"id": "1SPE-VARALEA-EXP-002", "chapitre": "1SPE-VARIABLES-ALEATOIRES", "type_objet": "experimentation", "capacites_codes": ["C6"], "programme_refs": ["VA-EXP-FONCTION-MOYENNE"], "mode_creation": "ex_nihilo", "sources_inspiration": [], "status": "verified"}
\subsection{Écrire une fonction renvoyant la moyenne d'un échantillon}

La taille $n$ d'un échantillon est le nombre de valeurs simulées. Si
$x_1,\ldots,x_n$ sont ces valeurs, leur moyenne est
\[
  m=\frac{x_1+\cdots+x_n}{n}.
\]

\begin{python}
from random import Random

VALEURS = (-2, 1, 4)
PROBABILITES = (0.2, 0.5, 0.3)

def tirage_variable(generateur):
    u = generateur.random()
    cumul = 0
    for valeur, probabilite in zip(VALEURS, PROBABILITES):
        cumul = cumul + probabilite
        if u < cumul:
            return valeur
    return VALEURS[-1]

def moyenne_echantillon(n, graine):
    """Renvoie la moyenne d'un echantillon simule de taille n."""
    if n <= 0:
        raise ValueError("n doit etre strictement positif")
    generateur = Random(graine)
    somme = sum(tirage_variable(generateur) for _ in range(n))
    return somme / n
\end{python}

\exemple{
Avec $n=40$ et la graine $1729$, la fonction renvoie $m=1{,}9$.
Cette valeur dépend de l'échantillon; elle n'est pas une nouvelle valeur
théorique de $E(X)$.
}

\subsection*{À toi de jouer — écrire la fonction}

On veut une fonction \code{ecart\_a\_esperance(n, graine)} qui renvoie
$\lvert m-E(X)\rvert$, où $m$ est la moyenne d'un échantillon de taille $n$.

\begin{enumerate}
  \item Recopier et compléter les deux lignes manquantes.
  \item Tester avec $n=40$ et la graine $1729$, puis avec $n=400$ et la même
        graine. Que constate-t-on ?
\end{enumerate}

\begin{python}
ESPERANCE = 1.3

def ecart_a_esperance(n, graine):
    moyenne = ...          # a completer : reutiliser moyenne_echantillon
    return ...             # a completer : la distance a ESPERANCE
\end{python}

\ifnxVersionProfesseur
\begin{corrige}{1SPE-VARALEA-EXP-002-ATOI}
\textbf{1.} \code{moyenne = moyenne\_echantillon(n, graine)} puis
\code{return abs(moyenne - ESPERANCE)}.

\textbf{2.} Avec $n=40$ on obtient $\lvert 1{,}9-1{,}3\rvert = 0{,}6$; avec
$n=400$ la moyenne vaut $1{,}54$ et l'écart tombe à $0{,}24$. Augmenter la taille de l'échantillon
rapproche la moyenne observée de l'espérance, sans jamais garantir l'égalité.
\end{corrige}
\fi

% META: {"id": "1SPE-VARALEA-EXP-003", "chapitre": "1SPE-VARIABLES-ALEATOIRES", "type_objet": "experimentation", "capacites_codes": ["C4"], "programme_refs": ["VA-EXP-DISTANCE-MOYENNE-ESPERANCE"], "mode_creation": "ex_nihilo", "sources_inspiration": [], "status": "verified"}

\subsection{Étudier la distance entre moyenne et espérance}

Pour la loi étudiée,
\[
  \mu=E(X)=1{,}3\qquad\text{et}\qquad \sigma=2{,}1.
\]
On mesure l'écart entre la moyenne $m$ d'un échantillon de taille $n$ et
l'espérance $\mu$ par la distance $\lvert m-\mu\rvert$.

\begin{python}
from random import Random

VALEURS = (-2, 1, 4)
PROBABILITES = (0.2, 0.5, 0.3)
MU = 1.3

def tirage_variable(generateur):
    u = generateur.random()
    cumul = 0
    for valeur, probabilite in zip(VALEURS, PROBABILITES):
        cumul = cumul + probabilite
        if u < cumul:
            return valeur
    return VALEURS[-1]

def distance_moyenne_esperance(n, graine):
    """Renvoie |m - mu| pour un echantillon simule de taille n."""
    if n <= 0:
        raise ValueError("n doit etre strictement positif")
    generateur = Random(graine)
    moyenne = sum(tirage_variable(generateur) for _ in range(n)) / n
    return abs(moyenne - MU)
\end{python}

\exemple{
Avec la graine $1729$, les distances obtenues pour
$n=20$, $n=100$ et $n=500$ valent respectivement environ
$0{,}30$, $0{,}18$ et $0{,}15$.
}

\erreurFrequente{
Dans cet exemple, la distance diminue quand $n$ augmente. Cette observation
empirique ne constitue pas une preuve qu'elle décroît à chaque nouvelle
simulation : un échantillon aléatoire peut exceptionnellement s'éloigner
de $\mu$.
}

% META: {"id": "1SPE-VARALEA-EXP-004", "chapitre": "1SPE-VARIABLES-ALEATOIRES", "type_objet": "experimentation", "capacites_codes": ["C7"], "programme_refs": ["VA-EXP-PROPORTION-2SIGMA"], "mode_creation": "ex_nihilo", "sources_inspiration": [], "status": "verified"}
\subsection{Simuler \texorpdfstring{$N$}{N} échantillons et tester la borne à \texorpdfstring{$2\sigma$}{2σ}}

On simule $N$ échantillons, chacun de taille $n$. Pour chaque échantillon,
on note $m$ sa moyenne et on teste
\[
  \lvert m-\mu\rvert\leq\frac{2\sigma}{\sqrt{n}}.
\]
Dans notre exemple, $\mu=1{,}3$ et $\sigma=2{,}1$. Les quatre paramètres
$\mu$, $\sigma$, $n$ et $N$ jouent des rôles différents : les deux premiers
décrivent la loi, $n$ est la taille d'un échantillon et $N$ est le nombre
d'échantillons simulés.

\begin{python}
from math import sqrt
from random import Random

VALEURS = (-2, 1, 4)
PROBABILITES = (0.2, 0.5, 0.3)
MU = 1.3
SIGMA = 2.1

def tirage_variable(generateur):
    u = generateur.random()
    cumul = 0
    for valeur, probabilite in zip(VALEURS, PROBABILITES):
        cumul = cumul + probabilite
        if u < cumul:
            return valeur
    return VALEURS[-1]

def proportion_dans_borne(n, N, graine):
    """Proportion des N moyennes a distance au plus 2*sigma/sqrt(n) de mu."""
    if n <= 0 or N <= 0:
        raise ValueError("n et N doivent etre strictement positifs")
    generateur = Random(graine)
    borne = 2 * SIGMA / sqrt(n)
    favorables = 0
    for _ in range(N):
        moyenne = sum(tirage_variable(generateur) for _ in range(n)) / n
        if abs(moyenne - MU) <= borne:
            favorables = favorables + 1
    return favorables / N
\end{python}

\exemple{
Pour $n=50$, $N=400$ et la graine $1729$, la proportion obtenue est
$0{,}94$. Avec la graine $314159$, elle vaut $0{,}935$ : la proportion
elle-même varie d'une simulation à l'autre.
}

\erreurFrequente{
Cette proportion est un constat empirique et ne constitue pas une preuve
d'une valeur universelle. Il faut conserver $n$, $N$, $\mu$, $\sigma$ et la
graine avec le résultat afin que l'expérience soit interprétable et
reproductible.
}

\subsection*{À toi de jouer — exploiter $N$ échantillons}

\begin{enumerate}
  \item Exécuter \code{proportion\_dans\_borne(50, 400, 1729)} puis reprendre
        avec la graine $20250901$. Les deux proportions sont-elles égales ?
        Pourquoi ?
  \item Reprendre avec $n=200$ à graine fixée. La borne
        $\frac{2\sigma}{\sqrt{n}}$ augmente-t-elle ou diminue-t-elle ?
        La proportion observée s'effondre-t-elle pour autant ?
  \item Rédiger en deux phrases ce que l'on peut, et ce que l'on ne peut pas,
        conclure d'une telle simulation.
\end{enumerate}

\ifnxVersionProfesseur
\begin{corrige}{1SPE-VARALEA-EXP-004-ATOI}
\textbf{1.} $0{,}94$ puis $0{,}9525$ : les proportions diffèrent car chaque
graine produit d'autres échantillons. L'ordre de grandeur, lui, se conserve.

\textbf{2.} Pour $n=200$ la borne passe de $0{,}594$ à $0{,}297$ : elle
diminue comme $1/\sqrt{n}$. La proportion reste voisine ($0{,}935$ à la graine
$1729$) car les moyennes se resserrent au même rythme que la borne.

\textbf{3.} On peut constater qu'une large majorité des moyennes tombe dans
l'intervalle annoncé, et que ce constat se reproduit sur d'autres graines. On
ne peut pas en déduire une valeur exacte ni une propriété démontrée : c'est
une observation empirique, dépendante de $n$, de $N$ et de la graine.
\end{corrige}
\fi

